\section{Introduction}
\label{sec:introduction}

Chain-of-thought reasoning makes intermediate solution steps available as explicit supervision~\citep{wei2022chain}.
By training on teacher-generated reasoning traces, smaller language models can learn from explanations in addition to final answers~\citep{hsieh2023distilling}.
The usefulness of this supervision, however, depends on what the traces contain.
A correct solution may include essential calculations alongside repeated explanations, alternative derivations, or verification steps that contribute little to a particular student's learning.
Longer traces therefore increase the amount of supervision without necessarily increasing its instructional value.
This raises a central question for reasoning distillation: how can we generate concise teaching traces that retain the information a smaller student needs?

Prior work has approached this question through the choice and construction of training data.
Small Models Struggle to Learn from Strong Reasoners shows that smaller students do not consistently benefit from longer reasoning chains and proposes mixing reasoning of different complexity~\citep{li2025small}.
LiteCoT uses difficulty-aware prompting to rewrite teacher traces into shorter solutions before distillation~\citep{wu2025concise}.
These findings motivate adapting reasoning supervision to the student, but they do not establish a universal advantage for shorter traces.
For example, Compress-Distill finds that compression reduces training and inference costs while original traces retain the highest accuracy in its main comparisons~\citep{griot2026compress}.
Together, these results suggest that length is an incomplete proxy for teaching quality: removing useful intermediate steps and removing redundant content can both shorten a trace, with different consequences for learning.

Length-based selection offers a simple way to obtain concise supervision, but it operates on completed samples rather than directly changing the teacher's generation process.
Activation steering provides a complementary approach.
ASC learns a steering direction from verbose--concise trace pairs and modifies internal activations to compress chain-of-thought generation~\citep{azizi2026activation}.
Steering vector distillation further demonstrates that behaviors induced in a teacher can transfer through training on its outputs~\citep{blank2026subliminal}.
These precedents motivate a more specific question: which features distinguish naturally occurring short and long correct solutions, and can manipulating those features produce useful supervision for a smaller student?
Answering this question requires connecting feature discovery, generation control, and downstream learning, since a reduction in teacher output length alone does not establish better distillation.

In this paper, we propose \textbf{SAE-guided reasoning distillation}, a feature-level intervention framework for producing concise teaching traces.
We begin with short and long solutions sampled from the same teacher for the same problems, retaining traces whose final answers are verified as correct.
This within-problem comparison reduces differences arising from teacher identity, problem content, and final-answer correctness.
We train a shared sparse autoencoder (SAE) on teacher activations from both groups, using a common dictionary to compare their feature activity~\citep{cunningham2023sparse,gao2024scaling}.
We then identify short-associated and long-associated features from their contrasting activation patterns.
These associations provide candidate intervention targets; they do not, by themselves, establish that an individual feature represents redundancy or a particular reasoning operation.

During generation, the framework uses the corresponding SAE decoder directions to encourage short-associated features and suppress long-associated features.
The resulting intervention acts while the teacher constructs a solution, with the aim of changing the distribution of teaching traces toward concise reasoning.
Complete, answer-verified outputs then provide supervision for student fine-tuning.
The SAE decomposition makes it possible to select and inspect individual components of the intervention, while downstream evaluation tests whether the generated data is useful for learning.
This separates the objective of controlling teacher verbosity from the objective of transferring reasoning capability.

Our empirical focus is GSM8K~\citep{cobbe2021training}, with a Qwen2.5-7B-Instruct teacher and a Qwen2.5-1.5B-Instruct student.
The evaluation distinguishes teacher trace length and correctness from student accuracy and output length, and separates equal-example from equal-supervision-token training budgets.
The central comparison is between students trained on intervention-generated traces and those trained on unmodified teacher outputs under controlled training conditions.
Naturally short traces and matched random interventions provide additional controls for selection effects and nonspecific activation perturbations.
This design tests whether feature-guided generation can improve the accuracy--efficiency trade-off of reasoning distillation, without assuming that shorter supervision is inherently better.

Our contributions are threefold: a within-problem approach to identifying length-associated teacher features using a shared SAE; a generation-time intervention framework that turns these features into candidate teaching data; and a controlled evaluation of how the resulting traces affect student learning and reasoning length.
Together, these components connect the internal representation of teacher verbosity to the construction and evaluation of reasoning supervision.
